% =============================================================================
%  HIKET -- MATERIALS AND METHODS, complete, for review (2026-10-21)
%
%  Standalone: compiles on its own (latexmk -pdf HIKET_methods_v2.tex) with the
%  same preamble as the drafts, and is meant to be merged as the Methods of
%  HIKET_draft_v2.tex once reviewed.
%
%  ORIGIN. The Methods of HIKET_draft_v1.tex (rounds 2-7 of Lorenzo's notes,
%  checked against the code), brought up to the storyline v4 design:
%    NEW      the pre-run climate (1961-1984 mean) in the transient initialisation
%    NEW      the comparison with a balanced start (the eqinit calibrations)
%    NEW      the division of the sink (one more calibration; recalibrated version)
%    NEW      forward runs: the projection under both starts, the input that would
%             hold the sink, and the warming scenarios with the published arm (F15)
%    CHANGED  order: setup -> calibration -> experiments built on the calibration
%             -> evaluation; software versions (R 4.6.1, BayesianTools 0.1.9, read
%             from the production metadata and the Roihu r-env module)
%  Every NEW or CHANGED passage is marked in a source comment; the round comments
%  of v1 were not carried over (they stay in v1 in git).
%
%  Runs the text describes: production 20261012_1042*, balanced start
%  20261012_12*-13*, balanced start with the climate held 20261013_13*.
%  STYLE: no imported jargon; no sentence-initial bare This/It/These; redundant
%  rather than compact; never refer to our own development history.
% =============================================================================
\documentclass[11pt]{article}
\usepackage[utf8]{inputenc}\usepackage[T1]{fontenc}
\usepackage[margin=2.5cm]{geometry}
\usepackage{graphicx}\usepackage{amsmath}\usepackage{amssymb}
\usepackage[hidelinks]{hyperref}
\usepackage[round]{natbib}
\usepackage{caption}\captionsetup{font=small,labelfont=bf}
\usepackage{xcolor}
\usepackage{titlesec}
\titlespacing*{\paragraph}{0pt}{0.8\baselineskip}{0.4em}
\linespread{1.08}

% marks a place where the text is not finished; nothing else should be coloured
\newcommand{\gap}[1]{\textcolor{red!70!black}{\small[\,#1\,]}}
\graphicspath{{figures/}}
\newcommand{\hfig}[4]{%   #1 file  #2 width  #3 label  #4 caption
  \begin{figure}[htbp]\centering
    \includegraphics[width=#2\linewidth]{#1}
    \caption{#4}\label{fig:#3}
  \end{figure}}
\newcommand{\suppfig}[1]{\textcolor{blue!55!black}{\small[Supplementary figure: \textsf{#1}]}}
\newcommand{\figref}[1]{\textcolor{blue!55!black}{\small[Figure: \textsf{#1}]}}
\newcommand{\sinput}{\sigma_{\mathrm{input}}}
\newcommand{\sinit}{\sigma_{\mathrm{init}}}

\title{\vspace{-1.2cm}HIKET --- Materials and methods\\[2pt]
\large complete version for review, to be merged into the second draft}
\author{}
\date{Version of \today}

\begin{document}
\maketitle

\section{Materials and methods}

% CHANGED 2026-10-21: one orienting paragraph, so the reader knows the order.
The section first describes the data (the soil carbon observations, the litter input and
the climate), then the six models and how they are started and calibrated, and then the
calibrations and simulations built on that common setup: the comparison with a soil
started in balance, the division of the modelled sink among its sources, and the
simulations beyond the observations. The evaluation of the models against the
observations closes the section.

% =============================================================================
\subsection{Study plots and soil carbon observations}\label{sec:plots}

The study is based on data from permanent forest plots of the Finnish National Forest
Inventory for which soil carbon has been measured in three campaigns: the eighth national
forest inventory (VMI8), the BioSoil campaign of 2006, and the Komeetta campaign of 2024.
After the exclusions described below, $456$ plots are available for calibration, carrying
$1205$ plot-year observations of soil carbon in total. Of the $456$ plots, $365$ plots
($80$ per cent, carrying $959$ observations) were used to fit the models, and $91$ plots
($20$ per cent, carrying $246$ observations) were held out and used only for evaluation.
Plots were assigned to the two groups at random, and a plot was assigned as a whole, so
that no plot contributes observations to both groups.

The first campaign is not a single year. Its plots were sampled region by region between
1986 and 1995: 69 plots in 1986, 76 in 1987, 103 in 1988, 54 in 1989 and 72 in 1995.
Averaged over the plots, the sampling year of the first campaign is 1989, and that is the
year at which the campaign is drawn whenever it is shown as a single point. In the
calibration, each observation is compared with the model in the year in which that plot
was actually sampled.

On the plots measured in all three campaigns ($310$ plots, the set on which every observed
rate of change in this paper is computed), the mean stock of the groups sampled in 1986,
1987, 1988, 1989 and 1995 differs by $18.3$ tonnes of carbon per hectare at first visit,
which might reflect a trend within that same campaign. The five groups of plots defined by
their first-campaign sampling year differ by $37.2$ tonnes in 2006 and by $34.2$ tonnes in
2024, years in which every plot of every group was measured in the same campaign. The
difference between the five groups is therefore a property of which plots each group
contains, and not of the year in which it was sampled, because the difference is larger
when the sampling year is held constant.

\gap{The exclusion rules (zero litter, absent organic layer, peat soils, constant litter
series, whole-profile stock above 250 t\,C\,ha$^{-1}$) are referred to as ``described below''
but are not yet written out; add one short paragraph listing them with the number of plots
each removes.}

% =============================================================================
\subsection{Soil carbon stocks on a single homogenised basis}\label{sec:soc}

Since the three campaigns were measured and processed under protocols that differed in
ways that matter for a comparison between them, all three campaigns were rebuilt on a
uniform basis, starting from the layer-by-layer carbon stocks measured in each of the three
campaigns and applying one set of conventions to all three. Those stocks are corrected for
the volume occupied by stones and for the bulk density of the fine earth. The volume of
stones was measured in 2006 on every plot by the rod-penetration method of
\citet{Viro1952}, in which a steel rod is driven into the mineral soil at the sampling
points and the mean depth it reaches is converted into a stone volume, and the same stone
volume is applied to the first campaign \citep{Lehtonen2016,Kramarenko}. The bulk density of
the mineral layers was predicted from the equations of \citet{TamminenStarr1994} for all
layers of the first campaign and for most layers of 2006, and measured directly for part of
the 2006 samples from the upper $20$ cm \citep{Kramarenko}. \gap{Confirm that the 2024
campaign uses the same stone volumes and the same bulk density equations.} The
coarse-fragment correction matters because the coarse-fragment content of these plots has a
median of $43$ per cent by volume, so a mineral stock computed without it is too large by a
factor of about $1.6$.

The quantity used as the calibration target is the whole soil profile: the measured organic
layer, plus the measured mineral soil to $40$ cm, plus a modelled deep tail, integrated only
as far down as the auger actually reached on that plot and in no case deeper than one metre.
The depth to which the profile is integrated is therefore a per-plot quantity and not a fixed
one metre, since using a fixed depth of, for example, one metre would invent soil carbon
below bedrock on the $38$ plots where augering stopped at $10$, $20$ or $40$ cm. The modelled
deep tail was validated against the measured $40$--$80$ cm layer of the BioSoil campaign,
which was held out of the fit: the median ratio of predicted to observed carbon in that layer
is $0.96$, with a bias of $-2.5$ tonnes of carbon per hectare.

The rebuilt target was checked against the national figures published for the two campaigns
for which such figures exist. Applying the same plot exclusions as the official calculation,
and integrating over the same layers, the rebuilt target returns weighted national means of
$59.06$ and $61.05$ tonnes of carbon per hectare for 2006 and 2024 against official values of
$59.06$ and $61.05$, that is, agreement to the nearest $0.01$ tonnes per hectare on $446$
plots.

Two further steps were taken with the first campaign, both of them established from the
structure of the measurements themselves rather than from the fit. First, the layers of the
three campaigns were standardised, because they were recorded differently: the organic layer
of the first campaign contains the fermentation and humus horizons only, whereas the organic
layer of the two later campaigns also contains the litter and moss layer. Each plot's own
litter and moss layer from 2006 was therefore added to its first-campaign organic layer, so
that the same compartment is compared across the three campaigns. Second, the first campaign
is treated as what it is, a campaign whose fieldwork ran over several years rather than a
campaign carried out in a single year. The year in which each plot was actually sampled is on
record, and those years, rather than one nominal year for the whole campaign, are the years
at which the models are compared with the observations.

On this final basis the median whole-profile stock is $62.0$ tonnes of carbon per hectare for
the first campaign, $66.4$ for 2006 and $67.2$ for 2024, and the observations vary between
plots with a coefficient of variation of $0.44$ to $0.53$ depending on the campaign.
\suppfig{F2\_baseline}

\gap{Data availability and acknowledgement: the soil carbon measurements are not a published
dataset, so the statement has to say where they come from, who compiled them and what of the
compiled form can be deposited. Settle this with the people who hold them before
submission.}

% =============================================================================
\subsection{Litter input}\label{sec:litter}

Litter input is taken from the annual reconstruction of tree biomass and litter input for the
permanent plots of the Finnish National Forest Inventory by \citet{TupekLitter}, published as
an open dataset. In that reconstruction, the tree biomass measured by the inventory is
interpolated to every year with the help of satellite-derived basal area, and litter is
computed as the biomass of each tree component multiplied by a turnover rate for that
component, plus the branches, foliage, stumps and roots left on site when a drop in basal
area identifies a harvest; the full method is given in \citet{TupekLitter}. The product
supplies, for each plot and year, the carbon in litter from trees, separated into non-woody
litter (foliage and fine roots), fine woody litter (branches and coarse roots) and coarse
woody litter (stumps), and each of the three classes is further divided into the four
chemical fractions used by the Yasso models, using proportions specific to each tree species
and component. The product does not contain a separate term for trees that die without being
harvested, and it does not include the understorey vegetation, because the product is a tree
litter product by design. Across the study plots the mean litter input carried by the product
is $\bar{J} = 2.511$ tonnes of carbon per hectare per year.

The national greenhouse gas inventory estimates the same flux for the mineral forest soils of
southern and northern Finland, from the forest inventory and biomass models
\citep{StatFin2025NID}. On the part that the two estimates have in common, the litter from
the turnover of living trees and from harvest residues, they agree in level in southern
Finland, where most of the study plots lie, and neither rises after about 2006; in northern
Finland the inventory estimate lies higher and keeps rising
\suppfig{S18\_litter\_vs\_inventory}. The inventory estimate also carries the litter of trees
that die without being harvested, a term that grows over the period and that the plot product
does not contain. \gap{confirm with LUKE that the two regions of the inventory data match the
region code of the study plots}

The first year of the product is treated separately. The median litter input of 1985 is
$0.060$ tonnes of carbon per hectare per year against $1.398$ in 1986, a step that the series
does not repeat in any later year. Because 1985 is both the first year of the simulation and
the year in which the first soil observations fall, the value for 1985 was reconstructed for
each plot and each chemical fraction by fitting a linear trend over 1986--1990 and
extrapolating one year backwards. The reconstruction places 1985 just below 1986, which is
the direction the growing-stock record requires.

% =============================================================================
\subsection{Climate}\label{sec:climate}

Climate forcing comes from the gridded daily weather observations of the Finnish
Meteorological Institute, interpolated from the station network onto a $10$ by $10$ km grid
by kriging with external drift \citep{Aalto2016} and extended by the Institute to the
present. The grids were obtained through the internal weather database of the Natural
Resources Institute Finland, which mirrors them from the Finnish Meteorological Institute.
From the grid cell nearest to each plot we took the daily mean, minimum and maximum
temperature, precipitation and global radiation for 1961--2025, and each model takes from
those series the variables it requires.
% NEW 2026-10-21: the two uses of the climate record (forward run vs pre-run).
The record is used in two ways. From 1985 onwards each year of the simulation is driven by
the climate observed in that year. Before 1985 the simulation is driven by one constant
climate, the mean of 1961--1984 at each plot (a plot mean of $2.2$ degrees Celsius, against
$2.7$ for 1985--2004), so that the soil enters 1985 adjusted to the climate it had lived in;
the same long-term mean serves as the reference climate wherever one is needed later
(Sections~\ref{sec:init}, \ref{sec:balanced} and \ref{sec:division}).

Because the response of decomposition to climate is not a straight line, the average of the
response over the years is not the response evaluated at the average climate. Averaging the
climate first and evaluating the response afterwards biases the climate modifier by $5$ to
$15$ per cent in boreal conditions. The climate modifier is therefore computed for each year
and averaged afterwards, in all six models.

% =============================================================================
\subsection{The six models}\label{sec:models}

Six models are used, and their roles are different by class. Yasso is the model this study is
focused on, because Yasso is the model the Finnish greenhouse gas inventory runs. All three
published versions of Yasso are included: Yasso07, Yasso15 and Yasso20. Alongside the three
versions of Yasso we included three simple models with one, two and three carbon pools,
denoted SP1, TP2 and TP3. The three simple models are there as a benchmark, since their
behaviour is transparent enough so that any result obtained with Yasso can be checked against
a structure simple enough to reason about quickly.

Due to their different structures and designs, the number of free parameters differs greatly
across the six models, from six free parameters in SP1 to twenty-six in Yasso15 and Yasso20.
The free parameters of the three versions of Yasso are the twelve fractions describing
transfers between pools, the climate parameters, the parameters describing the effect of the
size of woody litter, and the two auxiliary parameters described in
Section~\ref{sec:inputunc}.

The decomposition rates of Yasso are not fitted, and are held at their published values.
Holding the rates fixed is a deliberate choice to keep the model anchored to its original
definitions. Decomposition rates and transfer fractions do not act independently on the
carbon stock but are always interacting: a slower rate and a larger transfer into a slow pool
produce much the same trajectory, so fitting both together improves usually little and makes
the result harder to read.

Of the two, the rates are the more fundamental to the definition of each model, because the
rate of a pool is what distinguishes that pool from the others. Whether the pools of a soil
carbon model can also be identified with measurable fractions of organic matter is a
long-standing question in soil carbon modelling, and the answer has shaped how Yasso was
built. The first version of Yasso assumed that litter consists of groups of organic compounds,
each group decomposing at its own rate regardless of the litter it came from, and defined its
decomposition compartments accordingly as extractives, celluloses and lignin-like compounds,
followed by two compartments of humus \citep{Liski2005}. Yasso07 kept the same principle and
defined its four litter pools by fractions that can be measured on litter by sequential
extraction: the fraction soluble in acid, the fraction soluble in water, the fraction soluble
in ethanol, and the insoluble remainder \citep{Tuomi2009}. Defining the pools by the
chemistry of the litter is what makes Yasso operational, because the chemistry of litter can
be measured and supplied to the model directly. The same aim, of matching the pools of a
model with fractions that can be isolated in the laboratory, has been pursued for other
models, for example through fractionation schemes designed to correspond to the pools of
RothC \citep{Zimmermann2007}, and whether fractions isolated in the laboratory correspond to
the functional pools of models has been examined at length \citep{vonLutzow2007}.

Once carbon has entered the soil, the parameterisation of Yasso moves carbon between the four
chemical fractions in every direction. In the published parameterisation of Yasso07, for
example, $97$ per cent of the carbon decomposed from the insoluble fraction is passed to the
acid-soluble fraction, and $99$ per cent of the carbon decomposed from the acid-soluble
fraction is passed to the water-soluble fraction \citep{Tuomi2009}, which represents the
transformation of organic matter by decomposers rather than the decay of a fixed set of
compounds. The pools therefore enter the model defined by the chemistry of the litter, and
within the soil they act, as in any other compartmental soil carbon model, as compartments
defined by their rates of turnover.

The view of pools as compartments of turnover and the view of soil organic matter as a
continuum of quality are less opposed than they appear. A chemical fractionation is itself a
discretization: it divides a continuum into a small number of classes, the discretization is
made in the laboratory rather than in the model, and where the boundaries between classes fall
is set by the method chosen, so that the fractions are defined by the procedure that isolates
them \citep{vonLutzow2007}. The material being divided is not a set of discrete compounds.
Soil organic matter is the product of many processes of decomposition, microbial
transformation and recombination, much of that organic matter has no defined molecular
structure, and the chemical diversity of soil organic matter is large enough that the
diversity can be treated as a continuum \citep{LehmannKleber2015}.

The theory of continuous quality starts from that continuum: \citet{BosattaAgren1991}
describe soil organic matter by a distribution of quality, in which the quality of a unit of
carbon determines how fast that unit is decomposed, and decomposition in turn moves the carbon
along the continuum \citep{AgrenBosatta1998}. The same idea was formalised independently in
the geochemistry of sediments, where models with a few discrete fractions of different
reactivity are treated as discretizations of a continuous distribution of reactivity
\citep{BoudreauRuddick1991}. On this view, a pool of a compartmental model and a fraction of a
laboratory procedure are two discretizations of the same continuum, one made in the model and
one made in the laboratory, and in the model the rate of decomposition is what places each
pool on the continuum. Changing a rate therefore redefines the pool, whereas the transfer
fractions describe how carbon moves between pools that have already been defined. We
therefore keep the rates as the published parameterisation set them and calibrate the
transfer fractions. In all three versions of Yasso the decomposition rates, together with the
fraction of decomposed carbon that enters the humus pool, are fixed at their published
values: they are constants of the model and not parameters of the calibration, and no prior
is placed on them.

% =============================================================================
\subsection{The uncertainty of the litter input}\label{sec:inputunc}

The litter input is among the least certain of the quantities that drive a soil carbon model.
Litter is not measured on the plots but derived from tree biomass through biomass models and
turnover rates, parts of the input such as the litter of the understorey are missing from
tree-based products by construction, and the input carries year-to-year variation that
periodic inventories cannot capture. Uncertainty analyses of the Finnish soil carbon estimate
have identified the litter input, together with the initial state of the soil, as the dominant
sources of uncertainty \citep{Peltoniemi2006,LehtonenHeikkinen2015}. Those analyses propagated
the uncertainty forward, by sampling the input from assumed error distributions and running
the soil model with each sample. We instead estimate the error of the input from the soil
observations themselves, together with the parameters of the models, so that the correction
applied to the input and the uncertainty of that correction are both learned from the data
and carried into every result.

The correction is a single multiplier, $\sinput$, applied to the litter input of the product
over the whole simulation. The prior of $\sinput$ is log-normal with a centre of $1.08$, the
ratio of the total litter input to Finnish upland soils estimated by
\citet{LehtonenHeikkinen2015}, which includes the understorey, to the tree litter input of the
product used here, $2.70$ against $2.511$ tonnes of carbon per hectare per year, and a width
of $0.25$ on the logarithmic scale. Because the multiplier is one number for the whole
country, it corrects the national level of the input and not the input of individual plots,
whose errors are carried by the error model described below.
% CHANGED 2026-10-21: the arm-A sensitivity was run on an earlier setup only.
\gap{The sensitivity to the centre of this prior (centre $1.27$, the ratio implied by the
reconstruction of \citet{Liski2006}) was run on an earlier version of the setup and not
repeated on the present one: repeat it with the new input data, or drop the sentence.}

The same reasoning is extended to the starting state. The litter input of 1917, which sets the
state of the soil at the start of the simulation, is not known either, and is estimated
through a second parameter, $\sinit$, the ratio of the litter input of 1917 to the litter
input of 1985. The prior of $\sinit$ is log-normal with a centre of $0.90$, the ratio of the
growing stock of Finnish forests in 1917 to that in 1985 \citep{Korhonen2024} converted to a
ratio of litter input with the relation between litter input and growing stock observed in the
product, and the same width of $0.25$; the reconstruction of \citet{Liski2006}, based on the
same growing-stock record, implies a similar ratio, about $0.86$. Both parameters are estimated
together with the model parameters, so that the uncertainty of the input and of the starting
state propagates into the initial state, into the simulated trajectory and into the predictive
intervals of every model.

% =============================================================================
\subsection{Transient initialisation}\label{sec:init}

Rather than starting each plot from a soil in balance with its present litter input, each draw
of the parameter vector starts the soil in 1917 and runs it forward to the first soil
observation before the comparison with data begins. The year 1917 is a convention, chosen as
the year of Finnish independence and therefore the point from which the national forest record
is usually counted (any year starting far enough before the observations for the soil to carry
a history would work).
% CHANGED 2026-10-21: the pre-run climate stated.
The state of the soil in 1917 is obtained conventionally, with a steady state assumption with a
constant litter input equal to the litter input of 1917 and the constant climate of
1961--1984 (Section~\ref{sec:climate}). The litter input is then raised along a prescribed
shape over the sixty-eight years from 1917 to 1985, still under the same constant climate,
after which the simulation is driven by the annual litter and climate of the record. A
constant climate is used before 1985 because the record does not reach back to 1917; the
mean of the earliest observed decades is the closest available measure of the climate in
which the soil developed before the observations began.

The shape along which the litter input is raised is taken from the reconstruction of the
litter input to Finland's forest soils published by \citet{Liski2006}, who estimated that input
annually from 1922 onwards by combining national forest inventory data with biomass and
turnover models. Three decisions were made about that reconstruction. First, only the
tree-derived terms are used, meaning tree litter together with harvest residues and the litter
of trees that died naturally, because those are the terms closest to the litter product used
after 1985; that product represents natural mortality only when a death appears as a loss of
basal area, but natural mortality is between $2$ and $4$ per cent of the tree-derived input in
the reconstruction, and leaving it out changes the shape used here by less than $0.04$ in any
year. The ground vegetation is left out of the shape for the same reason that it is left out
after 1985, assuming that it is absorbed by the input multiplier described in
Section~\ref{sec:inputunc}, and including it before 1985 but not after would make that
multiplier mean two different things on either side of the join. Second, the national totals
are converted to a per-hectare basis by dividing by the area of upland forest soil.
\citet{Liski2006} report the carbon of upland soils both as a national total and per unit area,
and the ratio of the two gives an area of $13.90$ million hectares in 1922 and $15.22$ million
hectares in 2004. We placed the growth between the two values in 1965--1980, the period in
which \citet{Korhonen2024} locate most of the expansion of productive forest land. Third, only
the shape is taken and not the level: the series is rescaled to run from zero in 1917 to one in
1985, so that the reconstruction supplies only the path between the two endpoints. The two
endpoints are set by the two parameters described in Section~\ref{sec:inputunc}: the input at
the end of the path is the litter input of the product, averaged over 1985--1989 for each plot
and scaled by $\sinput$, and the input at the start of the path is a fraction $\sinit$ of the
input at the end.

\hfig{F4_initialization}{0.90}{init}{\textbf{How the starting state of the soil is obtained.}
Rather than starting each plot from a soil already in balance with its present litter input,
the simulation starts in 1917 from a soil in balance with the litter input of 1917 and the mean
climate of 1961--1984, and the litter input is then raised over the following sixty-eight years
along a shape taken from the reconstruction of Finland's forest litter input by
\citet{Liski2006}, after which the simulation is driven by the annual litter and climate of the
record. The litter input of 1917 is not known and is estimated as part of the calibration.}

% =============================================================================
\subsection{Keeping the models as comparable as possible}\label{sec:comparable}

The six models cannot be given the same external information. The models differ in structure
and in the number of their parameters, and the published literature provides different
information for each of them: the three versions of Yasso come with their own calibrations,
whereas the three simple models have no calibration for forest soils. The more the models
differ in the information they are given, in their assumptions and in their freedom, the less
a difference between their results can be attributed to their structure. We therefore reduced
that difference as far as the models allow, through one compromise for each class of
parameter, summarised in Table~\ref{tab:priors}. The compromises reduce the asymmetry between
the models without removing it, and what remains is stated with each compromise.

\subsubsection{Decomposition rates}

The rates of the three versions of Yasso are fixed at their published values, for the reasons
given in Section~\ref{sec:models}. The three simple models have no rates of their own for
forest soils, and take theirs from the ICBM model \citep{AndrenKatterer1997}: the fast rate is
fixed, and the slow rate or rates are given a very informative prior centred on the ICBM value,
so that the data can move them only slightly. The slow rates are pinned rather than fixed
because the ICBM values cannot be transferred to a forest soil unchanged. ICBM was calibrated
on an arable long-term experiment, and cropland and forest soils differ in what controls
decomposition: cultivation breaks up the soil aggregates that protect organic matter
physically in undisturbed soils \citep{Six2002}. Moreover, a rate is only meaningful together
with the climate response it was estimated with, and the simple models combine the rates of
ICBM with the climate response of Yasso07; the rates were rescaled once so that the two climate
responses agree at the climate of the experiment on which ICBM was calibrated, but the two
responses differ elsewhere. Fitting the rates freely, on the other hand, would give the simple
models a freedom that the versions of Yasso do not have, and a narrow prior around the ICBM
value is the compromise between the two.

\subsubsection{Transfer fractions}

The fractions that describe how carbon passes between pools are free in every model that has
them: twelve in each version of Yasso, and the fraction of decomposed carbon passed on to the
slower pool in the simple models, of which SP1 has none, TP2 one and TP3 two. All fractions
carry a prior of the same width, $0.4$ on the logit scale, centred on the published value of
each model, which for the simple models is the humification coefficient of ICBM.

\subsubsection{Climate response and size of woody litter}

The three versions of Yasso take the centres and widths of their climate priors from their own
published calibrations \citep{Tuomi2009,Tuomi2011woody}, and \gap{Yasso15/20 source: Yasso20 =
Viskari et al.\ 2022; Yasso15's calibration is unpublished (Järvenpää et al., in preparation)
-- decide whether to cite the FMI parameter files or the manuscript}. The simple models, which
have no published climate response of their own, borrow the climate response of Yasso07, centre
and width. Only the versions of Yasso describe how decomposition depends on the size of woody
litter, with priors taken from their own calibrations; the simple models have no such term.

\subsubsection{The two auxiliary parameters}

Every model carries the two parameters described in Section~\ref{sec:inputunc}, $\sinput$ and
$\sinit$, with the same prior, the same centre and the same width, in every model.

\begin{table}[htbp]\centering\small
\caption{\textbf{How the prior of each class of parameter was set.}\label{tab:priors}}
\gap{Parameter-class by prior-criteria table: for each class (decomposition rates, transfer
fractions, climate response, size of woody litter, the two auxiliary parameters), the source of
the centre, the source of the width, the transformation and constraint, and which models carry
it. The worked version is in \texttt{Prior\_specs/PRIOR\_HOMOGENIZATION\_PLAN.md}, section 2.}
\end{table}

% =============================================================================
\subsection{Numerical integration}\label{sec:integration}

Each model is run forward one year at a time, from 1917 to the last observation, with the
litter input and the climate of each year held constant within that year. Within a year, each
pool receives its share of the litter input and of the carbon passed on from other pools, and
loses a fixed fraction of its own carbon per unit of time. The carbon in the pools at the end of
the year can therefore be obtained in two ways. The simple way is to compute how fast each pool
is changing at the start of the year and to assume that the pool keeps changing at that speed
for the whole year. The exact way is to use the solution of the equations over the year, which
for a single pool is known in closed form: the pool approaches the value at which gains and
losses balance along an exponential curve, and never passes that value. No numerical solver of
differential equations is used: within each year the equations of all six models have a known
exact solution, and that solution is evaluated once per year. The only approximation is the
annual time step itself, and re-running the calibrated three-pool model with twelve monthly
steps a year changed the national mean stock by at most $0.7$ tonnes of carbon per hectare in
any year, about one per cent of the stock.

The two ways differ when a pool turns over faster than once a year. The simple way then carries
the pool past its balance value in one year and back in the next, and produces an oscillation
that belongs to the method of calculation and not to the model. The fastest pool of the
three-pool model turns over more than once a year in warm years, so the difference matters
here.

All six models use the exact solution. For several pools that pass carbon to one another the
exact solution involves the exponential of the matrix of rates and transfers, and the models
differ only in how that exponential is obtained. For the one-pool and two-pool models the
solution is written out as a formula. For the three-pool model carbon moves in one direction
only, from the fast pool to the slow pool and from the slow pool to the humus, so the solution
can also be written out as a formula, as a combination of one exponential decay per pool. In
the three versions of Yasso carbon moves between pools in every direction, and the exponential
is computed numerically in the Fortran code of the models by the standard method of scaling and
squaring \citep{MolerVanLoan2003}.

% =============================================================================
\subsection{Mean transit time as a common measure of a model}\label{sec:mtt}

The six models differ in how many carbon pools they contain and in how carbon moves between
those pools, so their parameters cannot be compared with one another directly. A single number
can be compared, however, and that number is the mean time a unit of carbon spends in the soil
between entering it as litter and leaving it as respired carbon dioxide. \citet{Sierra2017}
call that quantity the \emph{mean transit time}, and we use their term throughout. Much of the
soil carbon literature, including earlier work of our own, calls the same quantity the mean
residence time; Sierra and coauthors show that the term residence time has been used for at
least three different quantities and discourage its use, so we adopted their convention. At
steady state the mean transit time is also equal to the total carbon stock divided by the total
input flux, which is what many authors call the turnover time.

For a system in which each pool loses carbon in proportion to the amount of carbon it holds, and
in which the rates do not change with time, the mean transit time can be computed from the rates
and transfers of the model alone \citep{Rasmussen2016,Sierra2017}: at steady state the mean
transit time equals the total carbon the system holds divided by the flux of carbon entering it.
We therefore give each model a litter input of exactly one tonne of carbon per hectare per year
and compute the carbon stock at which the gains and losses of every pool balance. The balance
stock is computed directly and not by running the model until it settles: for the three simple
models the balance stock is written out as a formula, and for the three versions of Yasso the
balance equations of all pools are solved together as one linear system. The mean transit time
is then the balance stock divided by the input. Because the models are linear, the result does
not depend on the amount of input chosen: doubling the input doubles the stock and leaves the
ratio unchanged. An input of one tonne only makes the stock and the mean transit time the same
number; the quantity reported is the ratio, in years.

Two properties make the quantity useful as a basis for comparing models. First, the mean transit
time computed in this way is a property of each model alone. Neither the litter input
multiplier, nor the estimated starting state, nor any soil carbon observation enters the
calculation, so two models can be compared even though they were fitted to the same data with
different numbers of parameters, and our calibrated parameters can be compared with the
published parameters of the same model on exactly the same footing. For the same reason, the
mean transit time is not the ratio of simulated stock to simulated input at any date of the
simulation, which reflects the imbalance of the soil as well as the model. Second, the
calculation collapses a parameter vector of between six and twenty-six numbers into one number
with a physical meaning and a unit of years, which allows models that are otherwise very
different to be compared on the same axis.

Because a decomposition model responds to climate and to litter chemistry, the mean transit
time is not a constant of the model but a property of the model evaluated under a specific
condition for these variables. We therefore fix one reference condition and evaluate every
model, and every parameter set, under that same condition. The reference climate is the mean
climate of the study plots, with a mean annual temperature of $3.4$ degrees Celsius, a
temperature amplitude of $13.8$ degrees and an annual precipitation of $591$ millimetres. The
reference litter is the mean litter composition of the study plots, normalised to a total of
one, which divides the input into $0.766$ non-woody litter, $0.219$ fine woody litter and
$0.014$ coarse woody litter, each with its own chemical composition, so that the part of each
model describing the effect of the size of woody litter is exercised as it would be in a real
simulation. The mean transit times reported here should therefore be read as the mean transit
time of each model under such conditions, not as a universal constant of the model, nor as the
transit time of the real soil under changing conditions. The comparison between our calibration
and the published parameterisation does not depend on the choice of reference, because both
sides of the comparison are evaluated at the same reference.
% NOTE 2026-10-21: this reference climate (the plot mean over the observed years) is
% deliberately NOT the 1961-1984 reference of the initialisation; it describes the model,
% not the soil's history. Say so in the text if a reviewer may confuse the two.

% =============================================================================
\subsection{The error model}\label{sec:error}

Observations are compared with predictions on the logarithmic scale, so that the error is
proportional to the size of the stock. The total spread of the error is fixed at $0.800$ on the
logarithmic scale, about twice the spread of the observed stocks around the national mean of
each campaign, $0.39$. That spread is the error of an idealised model that reproduces exactly the
change of the national mean from one campaign to the next and knows nothing about individual
plots. Such a model would follow the measurements together with their errors, because the
measurements are all the information it has. A process-based model adds information that the
measurements do not contain, in the litter input, the climate and the representation of
decomposition, and where the measurements are in error a process-based model is expected to
depart from them. The error term has to leave room for that departure.

Four further reasons place the realistic error above that of the idealised model. The litter
input carries errors of its own at each plot: the variation of the litter input between plots,
$0.41$, is followed only weakly by the observed stocks, and counted as an independent error it
raises the total to $0.57$. The observed errors have heavier tails than a normal distribution.
The observed sample excludes outlying plots. Part of the measurement error, in the stone volume,
in the bulk density and in the modelled deep tail of the profile, is shared between plots and
does not appear in a spread between plots. A value at the upper end of the plausible range was
preferred, because a larger error weakens the pull of the observations relative to the priors,
so that results attributed to the observations are stated conservatively.

Observations of soil carbon are not independent of one another. Observations of the same plot
share whatever the models get systematically wrong about that plot, observations from the same
part of the country share whatever the models get wrong about that part of the country, and
observations from the same campaign share whatever the protocol of that campaign did
differently. The error of each observation is therefore divided into four parts: a part shared
by all observations in the same band of latitude, a part shared by all observations of the same
plot, a part shared by all observations of the same campaign, and a part that belongs to the
observation alone,
\[
  \log y_{ij} \;=\; \log f_{ij}(\theta) \;+\; u^{R}_{r(i)} \;+\; u^{P}_{i} \;+\; u^{C}_{j} \;+\; e_{ij},
\]
where $y_{ij}$ is the stock observed on plot $i$ in campaign $j$, $f_{ij}(\theta)$ the stock
predicted by the model, and $u^{R}$, $u^{P}$, $u^{C}$ and $e$ the four parts. Each part is
treated as a random error with a mean of zero and a fixed standard deviation, and no part is
estimated for any individual plot, band or campaign. Instead, the likelihood is averaged over
all the values each shared part could take, weighted by how probable each value is, which for
normal errors can be done exactly; the shared parts then disappear from the calculation and
leave only a correlation between observations that share a plot, a band or a campaign.

The standard deviations of the shared parts were taken from the observations themselves, after
removing the mean of each campaign. The standard deviation of the regional part, $0.117$, is the
spread between the means of eight bands of latitude, whose boundaries were placed so that each
band contains the same number of plots, because the plots are unevenly distributed from south to
north and bands of equal width would leave some bands almost empty. The standard deviation of the
plot part, $0.396$, is measured from how strongly two observations of the same plot vary
together. Independent errors add by their variances, that is by the squares of their standard
deviations, and not by the standard deviations themselves, so the variances of the four parts add
up to the variance of the total, $0.800^2$, and the standard deviation of the part that belongs
to each observation alone, $0.685$, is what remains once the three shared parts are accounted
for. Dividing the error in this way does not change how far a single observation may lie from the
model, but it changes how much the observations are worth together: observations that share part
of their error partly repeat one another, and the $1205$ observations carry about as much
information on the national mean as $220$ independent observations would. Both published
calibrations of Yasso assumed the errors of the observations to be independent
\citep{Tuomi2009,Viskari2022}, and the calibration of Yasso20 named the correlation of errors
within a time series as an effect still to be included \citep{Viskari2022}. Treating correlated
observations as independent leads to parameter uncertainties that are too narrow
\citep{Oberpriller2021}, and the shared parts of the error used here are a way of avoiding that.

The standard deviations of the campaign parts cannot be measured from the data, because a shared
offset between campaigns and a genuine change of the stock between campaigns look the same in the
observations. The standard deviations were therefore prescribed: $0.030$ for each of the two
later campaigns, which followed similar protocols, and twice that, $0.060$, for the first
campaign, which differs from the later two in three documented ways. The organic layer of the
first campaign excluded the litter and moss layer, which is added here from each plot's
measurement of 2006; the fieldwork of the first campaign extended over a decade; and the carbon
content of its mineral layers below $5$ cm was mostly predicted from the loss of mass on ignition
rather than measured directly \citep{Kramarenko}.

% =============================================================================
\subsection{Sampling}\label{sec:sampling}

% CHANGED 2026-10-21: versions read from the production metadata (R) and from the
% r-env module on the computing cluster (BayesianTools); v1 said 4.3.1 / 0.1.8.
Each model was calibrated with five independent runs of $50\,000$ iterations, of which the first
$5\,000$ were discarded, using the differential evolution Markov chain sampler of
\citet{terBraakVrugt2008}, in the implementation provided by the \texttt{BayesianTools} package
\citep{BayesianTools}, version $0.1.9$, under R $4.6.1$ \citep{RCore}. The sampler proposes each
new parameter vector from the difference between past states of its chains, occasionally along
the line joining the current state to a past one, so that the proposals adapt to the shape of the
posterior; this suits posteriors with strongly correlated parameters, such as the one described in
this paper. Each run maintains three interacting chains, so that the five runs together provide
fifteen chains. Every calibration in this paper, including those of
Sections~\ref{sec:balanced} and~\ref{sec:division}, uses the same settings.

The information that the observations add to each parameter was measured as the
Kullback--Leibler divergence of the marginal posterior from the marginal prior, estimated from
kernel densities of the samples. The divergence is zero when the posterior is identical to the
prior, that is when the observations say nothing about the parameter, and grows as the
observations narrow the distribution or move the distribution away from the prior. The divergence
is expressed in nats, the unit that results when information is measured with the natural
logarithm, as bits are the unit that results with the logarithm to base two. Two cases give a
sense of the scale: a posterior with the same centre as its prior and one quarter of its width
corresponds to about $0.9$ nats, and a posterior with the same width as its prior but moved by one
prior standard deviation corresponds to $0.5$ nats. Convergence was assessed with the potential
scale reduction factor and with the effective sample size, and the rate of proposals rejected as
numerically impossible was monitored throughout.
% NEW 2026-10-21: state the convergence of all 18 calibrations here once (checked on disk).
For all eighteen calibrations (six models, each calibrated with the transient initialisation,
with a balanced start, and with a balanced start and the climate held), the largest potential
scale reduction factor of any parameter was $1.009$ and the smallest effective sample size
$1716$.

% =============================================================================
\subsection{The comparison with a soil started in balance}\label{sec:balanced}

% NEW 2026-10-21. Code: Calibration_real_data_transient/equilibrium_init.R
% (HIKET_EQUILIBRIUM_INIT=1); runs 20261012_12*-13*; design and expectations in
% NEXT_RUN_equilibrium_init.md (written 2026-10-08, before the runs).
To test whether the transient initialisation is needed, every model was calibrated a second time
with the soil started in balance. The calibration is identical in every respect --- data, error
model, priors, sampling --- except for one factor: the starting state. In 1985 the soil is set to
the state at which its gains and losses balance under the litter input of 1985--1989, scaled by
$\sinput$, and the reference climate of 1961--1984; $\sinit$ is fixed at one, which makes the
litter path of Section~\ref{sec:init} flat, and is removed from the parameters. The prior of
$\sinput$ is the same as in the transient calibration. The soil then receives the observed litter
and climate after 1985, exactly as in the transient calibration.

The balanced start is the steady-state assumption applied at the start of the measurements, and
not the procedure of any particular inventory. The test is deliberately the strongest fair version
of that assumption: every other parameter is recalibrated, so the models are free to compensate
for the missing deficit through the litter input, the transfer fractions and the response to
climate. A soil started in balance in 1985 is not expected to stay flat, because the litter input
rose by about half until 2006; what distinguishes the two starts is the size of the accumulation
they produce, and whether a model started in balance can match the measured rates of
accumulation without an inherited deficit.

Four expectations were written down before the calibrations were run: that both starts would
reproduce the level of the campaign means; that the balanced start would produce too little
accumulation over 1985--2006 and 1985--2024 in most models; that it would compensate by moving
along the trade-off between litter input and transit time, with a larger $\sinput$, a shorter
transit time or both; and that Yasso07, whose single climate modifier rescales time in all pools
at once, would have the most room to compensate. The expectation that would have refuted the need
for the transient initialisation was also written down: that the balanced start reproduces the
observed rates within their intervals, at no material cost in fit, in most models.
\gap{Keep the expectations in the Methods (they make the test falsifiable) and report in the
Results which held and which did not (two did not: the balanced start overshoots the stock of the
first campaign, and its $\sinput$ went down rather than up)? Storyline v4 leaves the choice open.}

The two calibrations are compared on the trajectory of the national mean against the three
campaigns, on the rates of change over 1985--2006, 2006--2024 and 1985--2024, on the position of
each model along the trade-off between litter input and mean transit time, and on the fit. The
log-likelihoods of the two calibrations can be compared directly, because the data, the error
model and its spread are the same and the balanced start is a special case of the transient one;
the comparison is nevertheless made mainly on the rates, because the information criteria inherit
the overconfidence that the error model (Section~\ref{sec:error}) only partly removes.

% =============================================================================
\subsection{The division of the modelled sink}\label{sec:division}

% NEW 2026-10-21. Code: doublechecks/attribution_decomposition.R with
% HIKET_ATTR_MODE=recalibrated; runs: the two above plus 20261013_13*
% (HIKET_EQUILIBRIUM_INIT=1 + HIKET_CLIMATE_HELD=1). Full account in
% appendices/appendix_attribution.tex.
The sink each calibrated model produces over 1985--2024 was divided among three sources: the
state inherited from before 1985, the change of the litter input since 1985, and the climate. The
division is built in the same way as the comparison of Section~\ref{sec:balanced}, and on it. One
further calibration was made for every model, with the soil started in balance \emph{and} the
climate held at its reference, the mean of 1961--1984, in every year after 1985, in the
calibration as in the simulation. Each model is thereby calibrated three times on the same
measurements, adding one assumption at a time:
\begin{enumerate}\itemsep0pt
  \item the transient initialisation, with the observed litter and climate (the calibration used
    in every other result);
  \item a balanced start, with the observed litter and climate (Section~\ref{sec:balanced});
  \item a balanced start, with the observed litter and the climate held at the reference.
\end{enumerate}
The three parts of the sink are the differences between the three calibrations, year by year. The
inherited state is the sink of the first minus the sink of the second; the climate is the sink of
the second minus the sink of the third; and the litter change is the sink of the third, which
exists only because the litter departs from its 1985--1989 level, since a soil in balance under
an unchanging litter input and climate gains nothing. The three parts add up to the sink of the
first calibration by construction.

Because every calibration is fitted to the measurements under its own assumption, each part
compares two fitted models that differ by one assumption, and includes the change of the other
parameters that the data demand once that assumption is made. The division is therefore a
comparison between assumptions, and not the effect of one source on a model held fixed. The
climate changes the share of carbon that decomposes each year, so its effect depends on how much
carbon is present, and it has a joint effect with the inherited deficit and another with the
litter change. In the three parts the joint effects are assigned by the order of the
calibrations: the inherited part is measured under the observed climate and contains the joint
effect of climate and inheritance, and the climate part is measured on the observed litter and
contains the joint effect of climate and litter.
\gap{Five terms instead of three: the two joint effects can be reported as terms of their own with
two further calibrations (transient start with the climate held; balanced start with the litter
held at 1985--1989). Planned with the new input data; rewrite this paragraph then.}

The parts are computed for the plots measured in all three campaigns ($310$ plots) and for $50$
posterior draws of each calibration, and the national values are plain averages over plots and
draws, which keep the parts adding up to the sink. Two checks are made for every model: for every
draw of the two balanced-start calibrations, the soil with the litter and the climate both held at
their references keeps a sink of zero on every plot, and the three parts add up to the sink of the
first calibration in every year and draw. The size of the climate part depends on the choice of
the reference climate, whereas the ordering of the parts does not; the reference was chosen as the
climate the soil had lived in before the observations began, the same that drives the
initialisation.

% =============================================================================
\subsection{Simulations beyond the observations}\label{sec:forward}

% CHANGED 2026-10-21: extended from v1 (projection, equilibrium stock, +20%) with the
% projection under both starts, the input that holds the sink, and the warming
% scenarios with the published arm. Code: forward_scenarios.R (inside the predictive
% scripts), doublechecks/input_to_hold_sink.R, doublechecks/f15_forward_experiment.R,
% doublechecks/published_arm_dat.R; figures F15, F19.
Four analyses extend the calibrated models beyond the observations. They give orders of
magnitude for the accumulation that remains and for its sensitivity, and are not forecasts of the
national soil carbon balance: they describe the calibrated models on the study plots, per
hectare, under stated forcing.

\paragraph{The projection.} Each model is run for $60$ years after the last observation, with the
litter input held at its value in the last observed year, 2024, and with the climate of the last
$20$ observed years repeated three times, so that the projection contains neither a change of
climate nor a change of forest growth. Because the repeated climate alternates between a cooler
and a warmer decade, rates of change over the projection are computed over whole $20$-year cycles.
The projection uses $100$ posterior draws of each calibration and the plots measured in all three
campaigns, and is made both for the transient calibration and for the calibration started in
balance (Section~\ref{sec:balanced}), so that the effect of the starting assumption on the
projected sink can be read directly.

\paragraph{The distance from balance.} The equilibrium stock of each plot is the stock at which
gains and losses would balance if the litter input and climate of the projection were held
indefinitely, computed with the routine of each model for the state of balance, as for the mean
transit time (Section~\ref{sec:mtt}); the distance from balance is the equilibrium stock divided by
the stock of 2024, minus one.

\paragraph{More litter input.} The projection is repeated with the litter input increased by $20$
per cent from the first projected year onwards. Because the models are linear in their input, the
eventual gain of stock under that increase is $20$ per cent of the equilibrium stock, and the
scenario shows how much of the gain is realised within $60$ years. The same linearity gives, with
no further simulation, the increase of litter input that would hold a given sink over a given
period: for each posterior draw the sink over the period is a straight line in the input
multiplier, passing through the projections with the input held and increased by $20$ per cent,
and the multiplier that gives the target sink is read off that line. The targets are the observed
rates of 1985--2024 and 2006--2024, and the periods the first $10$ and $20$ years of the
projection. Because the projection carries no further warming, the multipliers obtained are lower
bounds for a warming climate.

\paragraph{Warming.} The response to warming is measured by adding a linear rise of temperature to
the repeated climate of the projection, from zero in 2025 to its full value in 2084, and comparing
the stock in 2084 with that of the projection without the rise. Three scenarios are taken from the
projections for Finland of \citet{RuosteenojaJylha2021}, SSP1-2.6, SSP2-4.5 and SSP5-8.5, whose
warming of the annual mean temperature is given relative to 1981--2010. Part of that warming has
already happened within the climate being repeated: on the study plots the mean of 2005--2024 is
$1.0$ degree warmer than that of 1985--2010, which stands in for 1981--2010. The full value of the
rise is therefore the warming still to come, $1.3$, $2.6$ and $5.4$ degrees for the three scenarios,
so that the warming already observed is not counted twice. The litter input is held at its 2024
value under warming, with no response of forest growth, so the rise isolates the response of
decomposition and biases every model toward a loss of carbon; the result is a sensitivity, not a
projection.

The response to warming is computed for three parameterisations of each version of Yasso: our
calibration ($40$ posterior draws), the calibration with a balanced start ($40$ posterior draws,
Section~\ref{sec:balanced}), and the published parameterisation. Both arms start from the
state built with the run's own pre-run climate (Section~\ref{sec:init}). For Yasso15 and Yasso20
the published arm uses draws from the published posterior distributions of the parameters; for
Yasso07, for which no such distribution is available, it uses the published parameter values. The
published kinetics are not refitted; only the litter input multiplier of the published arm is set,
for each draw, so that the arm reproduces the three campaign means, with $\sinit$ held at the median
of our calibration. Because the stock is exactly proportional to the litter input multiplier, the
adjustment sets the level without changing the kinetics, and the relative response to warming does
not depend on it.
The balanced-start calibration is run from its own starting state, the steady state under the
pre-run climate with $\sinit = 1$, and with its own posterior litter input multiplier; comparing its
response with that of our calibration shows whether relaxing the steady-state assumption changes
the sensitivity of the calibrated models to temperature, and not only their trajectory.

% =============================================================================
\subsection{Predictive evaluation and residual analysis}\label{sec:residuals_methods}

Predictions were generated for both the calibration plots and the held-out plots by drawing from
the posterior and running each draw forward. Model skill is reported as the coefficient of
determination and the root mean squared error, separately for the calibration plots and for the
held-out plots. Residuals were defined as the logarithm of the observed stock minus the logarithm
of the predicted stock, so that a positive residual means that the soil holds more carbon than the
model predicts. The residuals of the six models are nearly identical from plot to plot, with a
correlation of $0.99$ or higher between any two models, and the detailed analyses of the residuals
therefore use a single value for each observation, the mean of the residuals of Yasso07 and
Yasso15.

Three analyses were made. First, a random forest was fitted to the residuals of each model and
evaluated out of bag, with about sixty site variables describing the climate, the texture and
chemistry of the soil, the composition of the litter input and the stand in 1985, to measure how
much of the residual can be predicted at all and which variables lead. The ranking was checked with
a wider set of several hundred variables: the stand records and tree tallies of the four
inventories, the chemistry of the soil, topographic wetness, nitrogen deposition, the record of
management, and the cover of the understorey vegetation measured in 2023; variables that measure the
soil carbon itself were left out. The leading variables of the forest then guided two detailed
analyses, made on the mean residual of Yasso07 and Yasso15. Second, the mean residual of each plot
was regressed on the basal area of the stand in 1985, fitted as a smooth curve, together with three
properties of the site: the fertility of the site, expressed by the forest site type recorded by
the inventory; the texture of the mineral soil, recorded in the field in 1985; and the volume share
of stones in the soil. Third, the residual of each soil observation was related to the time elapsed
since the last regeneration cut before the soil was sampled, a regeneration cut being a felling
that removes the stand so that a new one can be established. The year of the cut was taken from the
stand records of the inventories of 1985, 1990, 1995 and 2021--2023 and from the forest use
declarations filed with the Finnish Forest Centre from 1997 onwards. Where no cut was recorded
before the sampling, the time since the stand was established was used instead, taken as the age of
the stand in 1985, measured on the tally trees, plus the years elapsed since 1985. Because each plot
contributes up to three observations, the comparison between the basal area of 1985 and the time
since the cut used a regression with a separate intercept for each plot.

\bibliographystyle{plainnat}
\bibliography{library}

\end{document}
